\documentclass[../DoAn.tex]{subfiles}
\begin{document}

\label{appendix:change-app}

\section{Tổng quan ứng dụng}

Ứng dụng \textit{demo-change-app} là hệ thống Quản lý Thay đổi và Go-Live (Change \& Go-Live Management) được xây dựng nội bộ nhằm mục đích kiểm chứng khả năng tích hợp của nền tảng IAM trong môi trường ứng dụng nghiệp vụ thực tế. Ứng dụng mô phỏng quy trình kiểm soát thay đổi phần mềm (change management) trong một tổ chức ngân hàng, trong đó mọi thay đổi lên môi trường sản xuất đều phải trải qua luồng phê duyệt có kiểm soát trước khi được triển khai.

Ứng dụng được xây dựng trên nền tảng Spring Boot (phía máy chủ) và Angular (phía giao diện), triển khai dưới dạng một container Docker duy nhất tại cổng 8085. Xác thực người dùng hoàn toàn được ủy thác cho \textit{iam-auth-service} thông qua luồng OAuth2 PKCE, do đó \textit{demo-change-app} không lưu trữ thông tin đăng nhập và không tự quản lý phiên làm việc. Mọi yêu cầu gửi đến phía máy chủ đều mang token JWT do IAM phát hành, và \textit{iam-gateway} thực thi kiểm tra quyền trước khi chuyển tiếp yêu cầu đến ứng dụng.

\section{Các thực thể nghiệp vụ}

Hệ thống \textit{demo-change-app} tổ chức dữ liệu xung quanh sáu thực thể nghiệp vụ chính. \textbf{Change Request} là thực thể trung tâm, đại diện cho một yêu cầu thay đổi phần mềm; mỗi yêu cầu bao gồm tiêu đề, mô tả, phạm vi ảnh hưởng, thời gian dự kiến triển khai và trạng thái hiện tại (nháp, chờ phê duyệt, đã phê duyệt, từ chối). \textbf{Approver} ghi nhận danh sách người phê duyệt được chỉ định cho từng yêu cầu thay đổi, cùng với quyết định và lý do tương ứng. \textbf{Go-Live Job} đại diện cho một lần thực thi triển khai (go-live) của một yêu cầu thay đổi đã được phê duyệt --- mỗi change request có thể tương ứng với đúng một go-live job, ghi nhận thời điểm thực thi, các bước triển khai theo thứ tự và trạng thái hoàn thành của từng bước. \textbf{Checklist Item} là các hạng mục kiểm tra bắt buộc phải hoàn tất trước, trong và sau khi một go-live job được xác nhận thành công. \textbf{Document} lưu trữ các tài liệu đính kèm liên quan đến yêu cầu thay đổi, bao gồm tài liệu kỹ thuật, kế hoạch rollback và biên bản phê duyệt. \textbf{Audit Log} ghi lại toàn bộ hành động của người dùng trong hệ thống, phục vụ mục đích truy vết và kiểm toán nội bộ.

Hệ thống định nghĩa bốn nhóm người dùng với tập quyền khác nhau: Developer có quyền tạo, chỉnh sửa và gửi yêu cầu thay đổi; Tech Lead có thêm quyền phê duyệt hoặc từ chối; Release Manager có quyền thực thi go-live job và cập nhật checklist; Admin nắm toàn bộ quyền bao gồm xem báo cáo và audit log.

\section{Luồng nghiệp vụ cơ bản}

Quy trình vận hành của \textit{demo-change-app} diễn ra qua hai giai đoạn liên tiếp. Trong giai đoạn phê duyệt, Developer tạo một change request, điền đầy đủ thông tin kỹ thuật và chỉ định danh sách người phê duyệt. Sau khi gửi chính thức, yêu cầu chuyển sang trạng thái chờ phê duyệt và hệ thống thông báo đến Tech Lead liên quan. Tech Lead xem xét nội dung, có thể yêu cầu bổ sung thông tin, rồi đưa ra quyết định chấp thuận hoặc từ chối kèm lý do. Chỉ khi tất cả người phê duyệt trong danh sách đồng ý, yêu cầu mới được chuyển sang trạng thái đã phê duyệt.

Trong giai đoạn triển khai, Release Manager nhận yêu cầu đã phê duyệt, khởi tạo go-live job và thực thi từng bước theo thứ tự đã định nghĩa. Trước khi xác nhận hoàn thành, toàn bộ checklist item phải được đánh dấu hoàn tất. Nếu phát sinh sự cố, Release Manager có thể kích hoạt rollback để đưa hệ thống về trạng thái trước khi thay đổi. Mọi hành động trong cả hai giai đoạn đều được ghi vào audit log tự động, không phụ thuộc vào thao tác thủ công của người dùng.

Ứng dụng này kiểm chứng rằng nền tảng IAM có khả năng kiểm soát quyền truy cập ở mức độ chi tiết: Developer không thể tự phê duyệt yêu cầu của chính mình, Tech Lead không thể thực thi go-live, và mọi ranh giới này được thực thi tập trung tại \textit{iam-gateway} mà không cần logic kiểm tra quyền phân tán trong mã nguồn ứng dụng.

\end{document}
